
% ============================================================
% SECTION 05
% INVARIANCE AND CONTINUITY
% Chapter 02: Invariance
% Book I: Ontology of Viability
% ============================================================

\section{Invariance and Continuity}
\label{sec:invariance-and-continuity}

% ------------------------------------------------------------
% 5.1. THE PROBLEM OF CONTINUITY
% ------------------------------------------------------------

The relationship between invariance and continuity
constitutes a further dimension of the ontological
problem developed in this chapter. Both concepts
concern the question of what remains identifiable
when an existent is considered under different
conditions. Nevertheless, they address distinct
questions and must not be treated as interchangeable.

Invariance, according to the provisional relational
definition adopted in this first edition, concerns
the relation between an existent, or one of its
determinate aspects, and a specified domain of
changes or transformations, in relation to which
something remains unchanged in a particular respect.

Continuity, by contrast, concerns whether and how
an existent continues through different conditions,
phases, or states. It raises a question about the
continuation of an existent rather than merely
the unchanged persistence of one of its determinations.

The distinction is important because the continuity
of an existent cannot be established simply by
identifying an invariant feature. Likewise, the
absence of an identified invariant feature does
not automatically establish the discontinuity
of an existent.

The present section examines this distinction
without presupposing a final account of persistence,
temporal identity, or the ontological ground of
continuity.

% ------------------------------------------------------------
% 5.2. CONTINUITY AS A DISTINCT QUESTION
% ------------------------------------------------------------

Continuity is often associated with the continuation
of something across different conditions or phases.
However, this association must be examined carefully.

To say that an existent continues is not necessarily
to say that every determination associated with it
remains unchanged. Different determinations may
undergo alteration while the question of the
continuity of the existent remains meaningful.

The concept of continuity therefore cannot be
reduced to the persistence of an unchanged
configuration or collection of features.

At the same time, the concept cannot be used
without qualification. A claim that an existent
continues requires some account of what is meant
by continuation and what distinguishes it from
the mere succession of different conditions.

The present analysis does not assume that
continuity is necessarily temporal in every
possible ontological context. It initially
examines continuity in relation to different
conditions or phases while leaving the broader
ontological scope of the concept open.

This restriction prevents the analysis from
imposing a temporal model on every form of
ontological continuity before the relevant
questions have been established.

% ------------------------------------------------------------
% 5.3. INVARIANCE DOES NOT ESTABLISH CONTINUITY
% ------------------------------------------------------------

The persistence of a particular determination
under a specified transformation does not,
by itself, establish the continuity of the
existent to which that determination belongs.

An invariant feature may be identified across
different conditions, but this identification
does not automatically explain whether the
existent itself continues through those
conditions.

The distinction can be expressed through
two separate questions:

\begin{enumerate}
    \item Does a particular determination
    remain unchanged under a specified
    domain of transformations?

    \item Does the existent continue through
    the conditions or phases being considered?
\end{enumerate}

The first question concerns invariance.
The second concerns continuity.

Although the answers may be related, neither
answer can be inferred from the other without
an additional argument.

For example, a particular structural relation
may remain unchanged across different
configurations. This may provide relevant
information about the transformations under
consideration, but it does not by itself
establish the ontological continuity of the
existent.

The relation between the invariant
determination and the continuity of the
existent requires further examination.

It would therefore be premature to define
continuity simply as the persistence of an
invariant feature.

% ------------------------------------------------------------
% 5.4. CONTINUITY DOES NOT REQUIRE THE
% INVARIANCE OF EVERY DETERMINATION
% ------------------------------------------------------------

If continuity were understood as the
unchanged persistence of every determination,
it would become difficult to account for
the possibility of continuation through
transformation.

An existent may undergo changes in
configuration, position, quantity, relation,
or other determinations. Such changes do
not automatically settle the question
of whether the existent continues.

The analysis must distinguish between
the continuation of the existent and
the persistence of each particular
determination associated with it.

This distinction does not establish that
an existent necessarily continues through
every possible transformation. It only
shows that the alteration of a particular
determination is not, by itself, sufficient
to establish discontinuity.

The conditions under which continuity
is maintained, altered, or interrupted
remain to be investigated.

Nor should the possibility of continuity
through change be used to introduce an
unexamined assumption about an immutable
substrate or an unchanging essence.

Such explanations may be considered in
comparative philosophical analysis, but
they cannot be adopted as foundational
principles without further justification.

% ------------------------------------------------------------
% 5.5. CONTINUITY AND THE SUCCESSION OF
% CONDITIONS
% ------------------------------------------------------------

A common way of examining continuity is
to consider an existent under successive
conditions.

For analytical purposes, these conditions
may be represented schematically as:

\[
A(t_1), \quad A(t_2), \quad A(t_3)
\]

The notation indicates the consideration
of an existent under different temporal
conditions. It does not, by itself,
establish that the same existent persists
through all of them.

The question of continuity remains distinct
from the formal representation of succession.

A sequence of descriptions may refer to
successive observations, successive states,
or successive conditions. The existence
of such a sequence does not automatically
establish the ontological continuity of
the subject being described.

The analysis must therefore distinguish
between:

\begin{enumerate}
    \item The succession of observations.
    \item The succession of described states.
    \item The continuation of an existent
    through different conditions.
\end{enumerate}

These are related but distinct matters.

The representation of temporal succession
may assist the analysis of continuity,
but it cannot substitute for an ontological
account of what it means for an existent
to continue.

The present chapter does not determine
whether time is the ultimate framework
of continuity or whether continuity
may have forms that cannot be reduced
to temporal succession.

% ------------------------------------------------------------
% 5.6. CONTINUITY, CHANGE, AND INVARIANCE
% ------------------------------------------------------------

The three concepts examined in this
chapter must now be distinguished
within a common analytical framework.

Change concerns the alteration of
one or more determinations under
specified conditions.

Invariance concerns the persistence
of a specified determination in
relation to a defined domain of
changes or transformations.

Continuity concerns whether and
how an existent continues through
different conditions or phases.

These concepts can be related without
being reduced to one another.

A possible case may involve change
in one determination, invariance in
another, and an unresolved question
about the continuity of the existent
as a whole.

Another case may involve continuity
without an immediately identifiable
invariant determination. Whether
such a case is ontologically possible
in a particular domain requires
further investigation.

The analytical distinctions can
be summarized as follows:

\[
\begin{aligned}
\text{Change} &\rightarrow
\text{Alteration of a determination} \\
\text{Invariance} &\rightarrow
\text{Unchanged determination within a domain} \\
\text{Continuity} &\rightarrow
\text{Continuation of an existent}
\end{aligned}
\]

These expressions are conceptual
summaries rather than exhaustive
definitions.

Their purpose is to prevent the
analysis from assuming that the
presence or absence of one of these
dimensions automatically determines
the status of the others.

% ------------------------------------------------------------
% 5.7. CONTINUITY AND IDENTITY
% ------------------------------------------------------------

The relationship between continuity
and identity is distinct from the
relationship between invariance
and identity.

The preceding chapter established
that being and identity are
ontologically inseparable. The
present section does not revise
that commitment.

However, the inseparability of being
and identity does not provide an
immediate explanation of how an
existent continues through different
conditions.

A claim of continuity concerns the
continuation of an existent, while
a claim about identity concerns
the identity-bearing character
of its existence.

These questions are related, but
they cannot be collapsed into a
single proposition.

It would be premature to conclude
that continuity is simply the
temporal persistence of identity,
because such a conclusion would
presuppose an account of the
relationship between identity
and temporal transformation.

It would also be premature to
conclude that continuity is
constituted by the persistence
of invariant features.

The relationship among identity,
continuity, and invariance must
therefore remain open to further
ontological development.

% ------------------------------------------------------------
% 5.8. EPISTEMIC RECOGNITION OF CONTINUITY
% ------------------------------------------------------------

The recognition of continuity
often involves the identification
of relations between different
conditions or phases.

A knower may compare descriptions,
observations, or records and
conclude that an existent has
continued through the conditions
being considered.

Such recognition depends upon
the available evidence, the
criteria of identification,
and the conceptual framework
within which the comparison
is conducted.

However, the epistemic recognition
of continuity must not be confused
with the ontological continuity
of the existent.

An existent may continue without
its continuity being recognized
by a particular knower. Conversely,
a mistaken identification may
lead an observer to attribute
continuity where the relevant
ontological conditions have not
been established.

The limitations of recognition
must therefore be distinguished
from the status of being itself.

Knowledge and non-knowledge remain
intertwined within the situated
condition of the knower. The
recognition of some relations
between successive conditions
may coexist with ignorance of
other relevant determinations.

This does not invalidate the
possibility of knowledge. It
requires that claims about
continuity remain proportionate
to the grounds on which they
are established.

% ------------------------------------------------------------
% 5.9. THE UNRESOLVED GROUND OF CONTINUITY
% ------------------------------------------------------------

The distinction between invariance
and continuity clarifies the
problem but does not resolve
its ontological ground.

Several questions remain open:

\begin{itemize}
    \item What makes the continuation
    of an existent possible?

    \item Is continuity an ontological
    determination, a relation between
    successive conditions, or a
    combination of distinct dimensions?

    \item Can continuity occur without
    the persistence of any particular
    determination identified as invariant?

    \item What is the relationship
    between continuity and the
    identity-bearing existence
    of an existent?

    \item Does continuity have the
    same meaning across different
    ontological realms?

    \item How does continuity relate
    to the dynamics of being?
\end{itemize}

The present chapter does not answer
these questions by presupposing
a permanent substance, an immutable
essence, a fixed structure, or
a universal temporal framework.

Such explanations would introduce
ontological commitments that have
not yet been established.

The question of continuity must
therefore remain open for the
subsequent development of the
system, particularly in relation
to the dynamics of being.

% ------------------------------------------------------------
% 5.10. SYNTHESIS AND TRANSITION
% ------------------------------------------------------------

The analysis has established that
invariance and continuity address
distinct but related ontological
problems.

Invariance concerns what remains
unchanged in a specified respect
within a defined domain of
transformations.

Continuity concerns whether and
how an existent continues through
different conditions or phases.

The persistence of an invariant
determination may be relevant
to the examination of continuity,
but it does not, by itself,
establish continuity.

Likewise, the continuation of
an existent does not necessarily
require the unchanged persistence
of every determination associated
with it.

The distinction prevents two
reductions:

\begin{enumerate}
    \item Continuity must not be
    reduced to the persistence
    of invariant features.

    \item Invariance must not be
    treated as a sufficient
    explanation of the continuation
    of an existent.
\end{enumerate}

The analysis also preserves the
foundational distinction between
identity and continuity. The
identity-bearing character of
existence does not, by itself,
settle the conditions under
which continuity is maintained
through transformation.

The next section examines the
relationship between invariance
and structure. This transition
is necessary because structural
relations may provide important
contexts in which invariance
can be identified, but their
persistence must not be
automatically interpreted as
the ultimate ground of either
identity or continuity.

% ============================================================
% END OF SECTION 05
% ============================================================
